\documentclass[aspectratio=169]{beamer}
\usepackage[utf8]{inputenc}
\usepackage[T1]{fontenc}

\title{A Comparative Analysis of Different Models for Bayesian Inference of Stratigraphic Intervals}
\author{Henrique L. Rieger \& Gustavo A. Ballen}
\institute{Instituto de Biociências -- Universidade Estadual Paulista ``Júlio de Mesquita Filho''}
\date{}


\begin{document}
  \maketitle

  \begin{frame}
    \frametitle{What are Stratigraphic Intervals?} \pause
    \begin{itemize}
      \item The range in which a species/lineage exists in time; \pause
      \item Can be measured in time or column depth; \pause
      \item Useful for: \pause
      \begin{itemize}
        \item Calibrating phylogenies; \pause
        \item Calculating evolutionary rates;
      \end{itemize}
    \end{itemize}
  \end{frame}

  \begin{frame}
    \frametitle{How do we determine a stratigraphic interval?} \pause
    \begin{itemize}
      \item Using the fossil record! \pause
      \item However, the fossil record is very incomplete; \pause
      \item But there is still loads of information there;
    \end{itemize}
  \end{frame}

  \begin{frame}
    \emph{*insert image of stratigraphic intervals of ammonites here*}

    \begin{itemize}
      \item $\theta_1$: Origin of lineage
      \item $\theta_2$: Extinction of lineage
    \end{itemize} \pause

    \begin{itemize}
      \item FAD: First Appearance Datum
      \item LAD: Last Appearance Datum
    \end{itemize} \pause

    The remainder fossils are useful to understand the underlying distribution.
  \end{frame}

  \begin{frame}
    \frametitle{Using Bayesian inference to estimate endpoints}

    \begin{equation}
      P(\theta\cdot|D) = \frac{P(\theta\cdot)P(D|\theta\cdot)}{\int_{\Theta}P(\theta)P(D|\theta)d\theta}
      \label{eq:bayes-theorem}
    \end{equation} \pause

    \begin{itemize}
      \item $P(\theta\cdot)$: prior distribution \pause
      \item $P(D|\theta\cdot)$: likelihood function \pause
      \item $P(\theta\cdot|D)$: posterior distribution \pause
    \end{itemize}

    Parameters of interest: $\theta_1$ and $\theta_2$; \pause

    We need some likelihood function that encodes endpoints as parameters.
  \end{frame}

  \begin{frame}
    \frametitle{The StratIntervals.jl package}

    Likelihood: Three-parameter Beta distribution (ref.)
    \begin{equation}
      f(\tau|\lambda,\theta_1,\theta_2) = 
        \begin{cases}
          \frac{(\theta_1 - \tau)^{-\lambda}}{(\theta_1 - \theta_2)^{1-\lambda}B(1, 1-\lambda)} &\text{if} \lambda \leq 0 \\
          \frac{(\tau - \theta_2)^{\lambda}}{(\theta_1 - \theta_2)^{1+\lambda}B(1+\lambda, 1)} &\text{if} \lambda > 0
        \end{cases}
      \label{eq:three-par-beta}
    \end{equation}

    Usage of MCMC algorithms provided by Turing (ref.);

    \emph{*insert GIF of three-par beta graph examples*}
  \end{frame}

  \begin{frame}
    \frametitle{Using four parameters}
    \begin{equation}
      f(\tau|\lambda,\theta_1,\theta_2) = \frac{(\theta_1 - \tau)^{\alpha-\beta}}{(\theta_1 - \theta_2)^{\alpha+\beta-1}B(\alpha, \beta)}
    \end{equation}

    More parameters = harder to converge \pause

    But greater variety of curves!

    \emph{*insert GIF of four-par beta graph examples*}
  \end{frame}
\end{document}
